% FILE: 06_contribution_to_thesis.tex
% Project: sources-pm4py — PM4Py Comparative Oracle Atlas
% Thesis role: pm4py-comparative-oracle-atlas

\section{Contribution to the Dissertation}
\label{sec:sources-pm4py-contribution}

\subsection{Contribution to Prediction vs.~Coordination}
\label{sec:sources-pm4py-contribution-prediction}

The oracle atlas makes a precise contribution to the dissertation's central argument about
prediction versus coordination. PM4Py's prediction/ML module (\texttt{pm4py.ml.*}) covers
next-activity prediction, remaining time prediction, outcome prediction, and decision-point
mining via decision trees --- and none of these are WASM-bridged (all \texttt{NO} in the
capability atlas). This is not a gap; it is a design boundary. The atlas establishes that
the wasm4pm bridge draws its boundary at the point where outputs can carry a deterministic,
named-law receipt. Predictive outputs cannot carry such receipts because they are
non-deterministic functions of model weights, which may not be traceable to a specific
event log under a specific algorithm and parameterisation. Coordination outputs --- fitness
scores, alignment costs, conformance diagnostics, soundness checks --- can carry receipts
because they are deterministic functions of a log, a model, and an algorithm. The oracle
atlas, by explicitly categorising the prediction capabilities as unbridged and explaining
why (via the NOT\_TO\_COPY manifesto's treatment of LLM integration and the divergence
matrix's treatment of LLM results as advisory-only), operationalises the thesis argument
that process coordination requires a different evidential standard than process prediction.
[AUTHOR THESIS / INTERPRETATION]

\subsection{Contribution to Process Evidence}
\label{sec:sources-pm4py-contribution-evidence}

The atlas's most direct contribution is to the dissertation's theory of process evidence.
It demonstrates, through the 37-algorithm crosswalk, that there is a systematic three-way
distinction between: type surfaces that must exist in wasm4pm-compat to carry receipts
(2 COVERED); type surfaces that exist for related structures but require additional work
to anchor receipts for specific algorithms (16 PARTIAL); and type surfaces that are
entirely absent, meaning no receipt is possible for the corresponding capability (19 MISSING).
This tripartite structure is the contribution: it says that process evidence is not a
binary property (a capability either exists or does not) but a graded property with a
type-law criterion for what constitutes admissible evidence.

The five JSON receipts provide the empirical grounding: they demonstrate that the
wasm4pm engine can produce fitness and precision scores that match PM4Py's deterministic
algorithms, with BLAKE3 provenance anchors and ed25519 signatures, for real enterprise
process claims (EBITDA rework reduction, working capital, SLA compliance). The
conformance-authority-map documents the mathematical foundations that underpin those
scores --- TBR formula $f(\sigma,N) = \tfrac{1}{2}(1-m/c)+\tfrac{1}{2}(1-r/p)$,
optimal alignment $\gamma^* = \operatorname{argmin}_\gamma \sum c(t,a)$, ETConformance
precision --- ensuring that the oracle's claims are traceable to published process mining
theory rather than to proprietary implementation choices.

\subsection{Contribution to Manufacturing Architecture}
\label{sec:sources-pm4py-contribution-manufacturing}

The atlas contributes to the dissertation's manufacturing architecture argument by
establishing the role of PM4Py as an oracle rather than a template. The NOT\_TO\_COPY
manifesto is the most direct formulation: ten architectural decisions that PM4Py made for
ergonomic reasons in the Python data science context would poison the manufacturing pipeline
if replicated in wasm4pm. DataFrame-first type erasure, mutable object construction,
generic exception raising, and silent projection loss are not simply inconveniences ---
they are structural incompatibilities with a manufacturing pipeline that requires every
artifact to have a provenance chain, every refusal to carry a named law type, and every
projection to carry a loss report.

The divergence matrix's operational checklist for porting PM4Py algorithms formalises
the manufacturing decision process: is the algorithm deterministic (must-match) or
heuristic (may diverge)? Does it produce a conformance artifact (must emit \texttt{Evidence<T,*,Witness>})?
Is it used downstream by other algorithms (output must be schema-locked)? Does it call
external services (results forbidden from conformance paths)? Does it perform lossy
conversion (must emit \texttt{LossReport})? This checklist is the manufacturing discipline
that distinguishes the research program's approach from ad hoc porting.

\subsection{Contribution to Capital-Flow Infrastructure}
\label{sec:sources-pm4py-contribution-capital}

The five JSON receipts connect the oracle atlas to the dissertation's capital-flow
infrastructure argument. Each receipt certifies a financial claim: the EBITDA rework
receipt asserts a \$1,250,000 annual EBITDA increase from Purchase Order rework reduction,
verified at fitness 0.982 and precision 0.945. The conformance-authority-map explicitly
addresses the role of conformance scores in M\&A diligence: fitness and precision are
not merely metrics --- they are the evidentiary basis for claiming that a process operates
within lawful bounds. The oracle atlas establishes that these claims are defensible only
when the conformance scores are produced by an engine that matches PM4Py's deterministic
algorithms exactly for the must-match capability set, with BLAKE3-anchored log and model
hashes and ed25519 signatures. Without the oracle atlas's mapping, there is no basis for
claiming that a wasm4pm fitness score is comparable to a PM4Py fitness score. With it,
the claim is traceable, law-bound, and receipt-anchored --- the precondition for
capital-flow defensibility. [AUTHOR THESIS / INTERPRETATION of receipt evidence]

\bigskip
\noindent\textit{Source authority:} CAPABILITY\_ATLAS.md (section 9: Prediction/ML, all
NO), NOT\_TO\_COPY.md (manufacturing architecture argument), ALGORITHM\_CROSSWALK.md
(Summary Table: 2 COVERED, 16 PARTIAL, 19 MISSING), divergence-matrix.md (section 13:
operational checklist), receipts/rec\_ebitda\_rework\_001.json (EBITDA claim and
verification results), conformance-authority-map.md (conformance formulas, M\&A
diligence context).
